\chapter{Conclusiones}

Se cumplieron los objetivos planteados al inicio del trabajo práctico, logrando adquirir destreza en el manejo de los controles fundamentales de un osciloscopio tanto analógico como digital.

A lo largo de las diversas experiencias, se comprendió la importancia de una correcta compensación de las puntas de prueba para evitar distorsiones en la medición de señales de alta frecuencia. Asimismo, se analizó el impacto del acoplamiento de entrada en la visualización de componentes continuas y alternas.

El uso de funciones avanzadas como el Hold-off, el disparo externo y el modo X-Y permitió resolver situaciones de medición complejas, como la estabilización de trenes de pulsos no periódicos o la comparación precisa de frecuencias.

Finalmente, la experiencia con el osciloscopio digital permitió apreciar las ventajas del almacenamiento y el barrido único para la captura de fenómenos transitorios, como el discado telefónico, facilitando el análisis detallado de eventos que serían difíciles de observar en un instrumento analógico convencional.
